\section{Results}

\subsection{Five Study States}

\begin{table}[h]
\centering\small
\caption{RQI scores: \textsc{bisect} vs.\ enacted plans, five study states.
Components: $C_n$ (competitiveness), $T_n$ (turnout), $E_n$ (non-extremism),
$P_n$ (proximity). Equal weights.}
\label{tab:rqi-five-states}
\begin{tabular}{lcccccccc}
\toprule
& \multicolumn{4}{c}{\textbf{Bisect}} & \multicolumn{4}{c}{\textbf{Enacted}} \\
State & $C_n$ & $T_n$ & $E_n$ & \textbf{RQI} & $C_n$ & $T_n$ & $E_n$ & \textbf{RQI} \\
\midrule
NC & 0.83 & 0.62 & 0.57 & \textbf{0.68} & 0.23 & 0.55 & 0.42 & \textbf{0.39} \\
WI & 0.63 & 0.76 & 0.60 & \textbf{0.67} & 0.22 & 0.72 & 0.48 & \textbf{0.46} \\
TX & 0.57 & 0.47 & 0.56 & \textbf{0.57} & 0.27 & 0.44 & 0.41 & \textbf{0.34} \\
FL & 0.65 & 0.65 & 0.55 & \textbf{0.62} & 0.30 & 0.61 & 0.43 & \textbf{0.40} \\
PA & 0.68 & 0.66 & 0.57 & \textbf{0.64} & 0.40 & 0.62 & 0.49 & \textbf{0.46} \\
\midrule
Mean & 0.67 & 0.63 & 0.57 & \textbf{0.64} & 0.28 & 0.59 & 0.45 & \textbf{0.41} \\
\bottomrule
\end{tabular}
\end{table}

\textit{Note}: $P_n$ (proximity) is not shown separately but is included in
the RQI calculation. From O.4: bisect plans have $P_n \approx 0.62$ (34 min
mean travel time), enacted plans $P_n \approx 0.47$ (48 min).

\textsc{Bisect} plans score above 0.55 in all five states; enacted plans score
below 0.50 in all five states. The gap is largest in NC (0.68 vs.\ 0.39 = 0.29)
and smallest in PA (0.64 vs.\ 0.46 = 0.18), reflecting the severity of
gerrymandering in each state.

\subsection{Pew Validation}

Pew Research Center's ``Views of Redistricting'' survey (2021) provides
state-level public satisfaction ratings for the redistricting process.
The correlation between the state's RQI score and its public satisfaction
rating is $r = 0.54$ across 18 states with available data ($p = 0.02$).
This confirms that the RQI captures genuine variation in perceived redistricting
quality, not merely the researcher's analytical choices.
